% =============================================================================
% State of the Art Presentation
% Explainable Disease Progression and Counterfactual Video Generation
% with Vision–Language Models
% =============================================================================

\documentclass[aspectratio=169]{beamer}
\usetheme{SimplePlus}

% -----------------------------------------------------------------------------
% Packages
% -----------------------------------------------------------------------------
\usepackage{subcaption}
\usepackage{hyperref}
\usepackage{amsmath,amssymb}
\usepackage{xcolor}
\usepackage{graphicx}
\usepackage{booktabs}
\usepackage{tikz}
\usepackage{multirow}

% -----------------------------------------------------------------------------
% Custom Colors
% -----------------------------------------------------------------------------
\definecolor{ForestGreen}{RGB}{34,139,34}
\definecolor{BrickRed}{RGB}{178,34,34}
\definecolor{DeepBlue}{RGB}{0,51,102}

% -----------------------------------------------------------------------------
% Title Information
% -----------------------------------------------------------------------------
\title[Explainable Disease Progression]{Explainable Disease Progression and Counterfactual Video Generation with Vision–Language Models}
\subtitle{State of the Art Review}
\author{Mohamed Amine}
\institute{TÉLUQ University, Montreal \\ Mitacs Globalink Research Internship \\ \vspace{0.5em} \small Supervisors: Dr. Belkacem Chikhaoui (Canada), Dr. Belkacem Khaldi (Algeria)}
\date{April 2026}

\begin{document}

% =============================================================================
% TITLE PAGE
% =============================================================================
\begin{frame}
    \titlepage
\end{frame}

% =============================================================================
% ROADMAP
% =============================================================================
\begin{frame}{Roadmap}
    \begin{columns}
        \begin{column}{0.5\textwidth}
            \begin{block}{Part I: Introduction}
                \begin{itemize}
                    \item Research Context \& Motivation
                    \item Problem Statement
                    \item Research Questions
                \end{itemize}
            \end{block}
            \begin{block}{Part II: Vision Transformers}
                \begin{itemize}
                    \item ViT Architecture
                    \item Swin Transformer
                    \item UNETR \& Swin UNETR
                \end{itemize}
            \end{block}
            \begin{block}{Part III: Language Models}
                \begin{itemize}
                    \item BioBERT \& PubMedBERT
                    \item Med-PaLM \& GPT-4 Medical
                    \item Vision-Language Integration
                \end{itemize}
            \end{block}
        \end{column}
        \begin{column}{0.5\textwidth}
            \begin{block}{Part IV: Generative Models}
                \begin{itemize}
                    \item Diffusion Models (DDPM)
                    \item Latent Diffusion \& Video LDM
                    \item Counterfactual Generation
                \end{itemize}
            \end{block}
            \begin{block}{Part V: Datasets}
                \begin{itemize}
                    \item BraTS Challenge
                    \item TCGA \& LIDC-IDRI
                    \item MIMIC-III Clinical Data
                \end{itemize}
            \end{block}
            \begin{block}{Part VI: Project Plan}
                \begin{itemize}
                    \item Research Objectives
                    \item Timeline \& Deliverables
                \end{itemize}
            \end{block}
        \end{column}
    \end{columns}
\end{frame}

% =============================================================================
% TRANSITION 1: Introduction
% =============================================================================
\begin{frame}[plain]
    \centering
    \vfill
    \begin{beamercolorbox}[center,shadow=true,rounded=true]{title}
        \usebeamerfont{title}
        \textbf{Part I: Introduction}\\
        \vspace{0.5em}
        \large Research Context, Motivation \& Problem Statement
    \end{beamercolorbox}
    \vfill
    
    \vspace{1em}
    \normalsize
    \textit{Why do we need explainable AI for disease progression modeling?}
\end{frame}

% =============================================================================
% INTRODUCTION FRAMES
% =============================================================================
\begin{frame}{Research Context: The Challenge}
    \begin{columns}
        \begin{column}{0.5\textwidth}
            \begin{block}{Current Medical Imaging Landscape}
                \begin{itemize}
                    \item \textbf{Data Explosion:} Medical imaging data growing exponentially
                    \item \textbf{Complexity:} Multi-modal, longitudinal patient data
                    \item \textbf{Critical Decisions:} Tumor monitoring, treatment planning
                    \item \textbf{Radiologist Burden:} Increasing workload, fatigue errors
                \end{itemize}
            \end{block}
        \end{column}
        \begin{column}{0.5\textwidth}
            \begin{block}{AI in Healthcare: The Promise}
                \begin{itemize}
                    \item Deep learning achieves expert-level performance
                    \item Automated lesion detection and segmentation
                    \item Quantitative disease tracking
                    \item \textbf{BUT:} Models remain ``black boxes''
                \end{itemize}
            \end{block}
            \begin{alertblock}{The Trust Gap}
                Clinicians hesitate to adopt AI systems they cannot understand or verify.
            \end{alertblock}
        \end{column}
    \end{columns}
\end{frame}

\begin{frame}{Problem Statement: Three Key Challenges}
    \begin{block}{Challenge 1: Explainability Crisis}
        Current deep learning models for medical imaging lack interpretability:
        \begin{itemize}
            \item \textbf{Why} did the model predict tumor growth?
            \item \textbf{What} features drove the diagnosis?
            \item \textbf{How} confident is the prediction?
        \end{itemize}
    \end{block}
    
    \begin{block}{Challenge 2: Temporal Modeling Gap}
        Disease progression is inherently temporal, yet most models are static:
        \begin{itemize}
            \item Single time-point analysis misses progression patterns
            \item Longitudinal data underutilized in current pipelines
            \item Need for video-like representation of disease evolution
        \end{itemize}
    \end{block}
    
    \begin{block}{Challenge 3: Counterfactual Reasoning}
        Clinicians need to answer ``what-if'' questions:
        \begin{itemize}
            \item ``What would happen without treatment?''
            \item ``How would the tumor evolve with alternative therapy?''
        \end{itemize}
    \end{block}
\end{frame}

\begin{frame}{Research Questions}
    \begin{alertblock}{Main Research Question}
        \centering
        \large How can we develop an explainable AI system that models disease progression \\
        and generates counterfactual visualizations using vision–language models?
    \end{alertblock}
    
    \vspace{1em}
    
    \begin{columns}
        \begin{column}{0.5\textwidth}
            \begin{block}{Sub-Questions}
                \begin{enumerate}
                    \item How can Vision Transformers extract clinically meaningful representations from medical images?
                    \item How can LLMs provide natural language explanations for visual features?
                    \item How can diffusion models generate realistic counterfactual disease trajectories?
                \end{enumerate}
            \end{block}
        \end{column}
        \begin{column}{0.5\textwidth}
            \begin{block}{Expected Contributions}
                \begin{itemize}
                    \item Novel ViT-LLM integration framework
                    \item Explainable disease progression model
                    \item Counterfactual video generation pipeline
                    \item Clinical evaluation methodology
                \end{itemize}
            \end{block}
        \end{column}
    \end{columns}
\end{frame}

% =============================================================================
% TRANSITION 2: Vision Transformers
% =============================================================================
\begin{frame}[plain]
    \centering
    \vfill
    \begin{beamercolorbox}[center,shadow=true,rounded=true]{title}
        \usebeamerfont{title}
        \textbf{Part II: Vision Transformers}\\
        \vspace{0.5em}
        \large From ViT to UNETR: Transformers for Medical Imaging
    \end{beamercolorbox}
    \vfill
    
    \vspace{1em}
    \normalsize
    \begin{block}{Why Vision Transformers?}
        \begin{itemize}
            \item Global context modeling through self-attention
            \item Natural alignment with language models
            \item State-of-the-art performance in medical image analysis
        \end{itemize}
    \end{block}
\end{frame}

% =============================================================================
% VISION TRANSFORMER FRAMES
% =============================================================================
\begin{frame}{Vision Transformer (ViT): The Foundation}
    \begin{columns}
        \begin{column}{0.5\textwidth}
            \begin{block}{Key Innovation (Dosovitskiy et al., 2020)}
                Apply pure Transformer architecture to images:
                \begin{enumerate}
                    \item Split image into fixed-size patches
                    \item Flatten patches into vectors
                    \item Add positional embeddings
                    \item Process with standard Transformer encoder
                \end{enumerate}
            \end{block}
            
            \begin{block}{Patch Embedding}
                \begin{center}
                    Image $\rightarrow$ Patches $\rightarrow$ Linear Projection $\rightarrow$ Tokens
                \end{center}
                For $224 \times 224$ image with $16 \times 16$ patches: \\
                $N = (224/16)^2 = 196$ tokens
            \end{block}
        \end{column}
        \begin{column}{0.5\textwidth}
            \begin{block}{Self-Attention Mechanism}
                \begin{center}
                    $\text{Attention}(Q, K, V) = \text{softmax}\left(\frac{QK^T}{\sqrt{d_k}}\right)V$
                \end{center}
                \begin{itemize}
                    \item $Q$, $K$, $V$: Query, Key, Value matrices
                    \item $d_k$: Dimension of keys
                    \item Global receptive field from layer 1
                \end{itemize}
            \end{block}
            
            \begin{alertblock}{Medical Imaging Advantage}
                Self-attention captures long-range dependencies critical for:
                \begin{itemize}
                    \item Tumor boundary delineation
                    \item Multi-organ relationships
                    \item Bilateral symmetry analysis
                \end{itemize}
            \end{alertblock}
        \end{column}
    \end{columns}
\end{frame}

\begin{frame}{ViT Architecture: Mathematical Formulation}
    \begin{block}{Patch Embedding Process}
        \begin{center}
            \Large
            $\mathbf{z}_0 = [\mathbf{x}_{\text{class}}; \mathbf{x}_p^1\mathbf{E}; \mathbf{x}_p^2\mathbf{E}; \cdots; \mathbf{x}_p^N\mathbf{E}] + \mathbf{E}_{\text{pos}}$
        \end{center}
        \begin{itemize}
            \item $\mathbf{x}_p^i \in \mathbb{R}^{P^2 \cdot C}$: Flattened patch $i$
            \item $\mathbf{E} \in \mathbb{R}^{(P^2 \cdot C) \times D}$: Patch embedding projection
            \item $\mathbf{E}_{\text{pos}} \in \mathbb{R}^{(N+1) \times D}$: Learnable positional embeddings
            \item $\mathbf{x}_{\text{class}}$: Classification token
        \end{itemize}
    \end{block}
    
    \begin{block}{Transformer Encoder Layers}
        \begin{center}
            $\mathbf{z}'_\ell = \text{MSA}(\text{LN}(\mathbf{z}_{\ell-1})) + \mathbf{z}_{\ell-1}$ \hspace{2em} (Multi-Head Self-Attention)\\
            \vspace{0.5em}
            $\mathbf{z}_\ell = \text{MLP}(\text{LN}(\mathbf{z}'_\ell)) + \mathbf{z}'_\ell$ \hspace{2em} (Feed-Forward Network)
        \end{center}
    \end{block}
    
    \begin{block}{Classification Output}
        \begin{center}
            $\mathbf{y} = \text{LN}(\mathbf{z}_L^0)$ \hspace{2em} (Use class token from final layer)
        \end{center}
    \end{block}
\end{frame}

\begin{frame}{Swin Transformer: Hierarchical Vision Transformer}
    \begin{columns}
        \begin{column}{0.5\textwidth}
            \begin{block}{Key Innovation (Liu et al., 2021)}
                \textbf{Shifted Window Attention:}
                \begin{itemize}
                    \item Local attention within windows
                    \item Shifted windows enable cross-window connections
                    \item Linear complexity: $\mathcal{O}(n)$ vs $\mathcal{O}(n^2)$
                \end{itemize}
            \end{block}
            
            \begin{block}{Hierarchical Feature Maps}
                \begin{itemize}
                    \item Stage 1: $\frac{H}{4} \times \frac{W}{4}$, $C$ channels
                    \item Stage 2: $\frac{H}{8} \times \frac{W}{8}$, $2C$ channels
                    \item Stage 3: $\frac{H}{16} \times \frac{W}{16}$, $4C$ channels
                    \item Stage 4: $\frac{H}{32} \times \frac{W}{32}$, $8C$ channels
                \end{itemize}
            \end{block}
        \end{column}
        \begin{column}{0.5\textwidth}
            \begin{block}{Window-Based Self-Attention}
                \begin{center}
                    $\text{Attention}(Q, K, V) = \text{softmax}\left(\frac{QK^T}{\sqrt{d}} + B\right)V$
                \end{center}
                \begin{itemize}
                    \item $B$: Relative position bias
                    \item Window size: $M \times M$ (typically $M=7$)
                    \item Complexity: $\mathcal{O}(M^2 \cdot n)$
                \end{itemize}
            \end{block}
            
            \begin{alertblock}{Medical Imaging Benefits}
                \begin{itemize}
                    \item Multi-scale feature extraction
                    \item Efficient for high-resolution images
                    \item Compatible with CNN decoders (U-Net)
                \end{itemize}
            \end{alertblock}
        \end{column}
    \end{columns}
\end{frame}

\begin{frame}{UNETR: Transformers Meet U-Net}
    \begin{columns}
        \begin{column}{0.5\textwidth}
            \begin{block}{Architecture Overview (Hatamizadeh et al., 2022)}
                \textbf{Encoder:} Pure ViT for feature extraction
                \begin{itemize}
                    \item 3D patch embedding for volumetric data
                    \item 12 Transformer blocks
                    \item Skip connections from layers 3, 6, 9, 12
                \end{itemize}
                
                \vspace{0.5em}
                \textbf{Decoder:} CNN-based upsampling
                \begin{itemize}
                    \item Progressive upsampling
                    \item Skip connections from encoder
                    \item Final segmentation head
                \end{itemize}
            \end{block}
        \end{column}
        \begin{column}{0.5\textwidth}
            \begin{block}{3D Patch Embedding}
                For 3D medical volumes ($H \times W \times D$):
                \begin{center}
                    $N = \frac{H}{P} \times \frac{W}{P} \times \frac{D}{P}$
                \end{center}
                With $P=16$, a $128^3$ volume yields $8^3 = 512$ tokens
            \end{block}
            
            \begin{block}{Performance on BraTS}
                \begin{center}
                    \begin{tabular}{lc}
                        \toprule
                        \textbf{Metric} & \textbf{Dice Score} \\
                        \midrule
                        Whole Tumor & 0.912 \\
                        Tumor Core & 0.857 \\
                        Enhancing Tumor & 0.823 \\
                        \bottomrule
                    \end{tabular}
                \end{center}
            \end{block}
        \end{column}
    \end{columns}
\end{frame}

\begin{frame}{Swin UNETR: State-of-the-Art for 3D Medical Segmentation}
    \begin{block}{Key Innovation (Tang et al., 2022)}
        Replace ViT encoder with Swin Transformer for improved efficiency and performance
    \end{block}
    
    \begin{columns}
        \begin{column}{0.5\textwidth}
            \begin{block}{Architecture Improvements}
                \begin{itemize}
                    \item \textbf{Hierarchical encoder:} Multi-scale features
                    \item \textbf{Shifted 3D windows:} Efficient attention
                    \item \textbf{Self-supervised pretraining:} Better initialization
                \end{itemize}
            \end{block}
            
            \begin{block}{Self-Supervised Learning}
                Pretrained on 5,050 CT volumes using:
                \begin{itemize}
                    \item Masked volume inpainting
                    \item Contrastive learning
                    \item Rotation prediction
                \end{itemize}
            \end{block}
        \end{column}
        \begin{column}{0.5\textwidth}
            \begin{block}{BraTS 2021 Results}
                \begin{center}
                    \begin{tabular}{lcc}
                        \toprule
                        \textbf{Region} & \textbf{UNETR} & \textbf{Swin UNETR} \\
                        \midrule
                        Whole Tumor & 0.912 & \textcolor{ForestGreen}{0.929} \\
                        Tumor Core & 0.857 & \textcolor{ForestGreen}{0.891} \\
                        Enhancing & 0.823 & \textcolor{ForestGreen}{0.854} \\
                        \bottomrule
                    \end{tabular}
                \end{center}
            \end{block}
            
            \begin{alertblock}{Project Relevance}
                \textbf{Swin UNETR} will be our primary architecture for:
                \begin{itemize}
                    \item Tumor segmentation
                    \item Feature extraction for progression modeling
                \end{itemize}
            \end{alertblock}
        \end{column}
    \end{columns}
\end{frame}

\begin{frame}{Vision Transformers: Summary Comparison}
    \begin{center}
        \begin{tabular}{@{}p{2.5cm}p{3cm}p{3.5cm}p{4cm}@{}}
            \toprule
            \textbf{Model} & \textbf{Attention Type} & \textbf{Complexity} & \textbf{Medical Application} \\
            \midrule
            ViT & Global self-attention & $\mathcal{O}(n^2)$ & Classification, feature extraction \\[8pt]
            Swin Transformer & Shifted window & $\mathcal{O}(n)$ & Multi-scale analysis, detection \\[8pt]
            UNETR & ViT + CNN decoder & $\mathcal{O}(n^2)$ & 3D segmentation \\[8pt]
            Swin UNETR & Swin + CNN decoder & $\mathcal{O}(n)$ & 3D segmentation (SOTA) \\
            \bottomrule
        \end{tabular}
    \end{center}
    
    \vspace{1em}
    
    \begin{alertblock}{Our Approach}
        \begin{itemize}
            \item Use \textbf{Swin UNETR} for segmentation and feature extraction
            \item Extract intermediate representations for disease progression modeling
            \item Connect ViT features to language models for explainability
        \end{itemize}
    \end{alertblock}
\end{frame}

% =============================================================================
% TRANSITION 3: Language Models
% =============================================================================
\begin{frame}[plain]
    \centering
    \vfill
    \begin{beamercolorbox}[center,shadow=true,rounded=true]{title}
        \usebeamerfont{title}
        \textbf{Part III: Large Language Models for Medical AI}\\
        \vspace{0.5em}
        \large From BioBERT to Med-PaLM: Natural Language Understanding in Healthcare
    \end{beamercolorbox}
    \vfill
    
    \vspace{1em}
    \normalsize
    \begin{block}{Why Language Models?}
        \begin{itemize}
            \item Generate natural language explanations for visual features
            \item Bridge the gap between AI predictions and clinical understanding
            \item Enable counterfactual reasoning through text descriptions
        \end{itemize}
    \end{block}
\end{frame}

% =============================================================================
% LANGUAGE MODEL FRAMES
% =============================================================================
\begin{frame}{BioBERT \& PubMedBERT: Domain-Specific Language Models}
    \begin{columns}
        \begin{column}{0.5\textwidth}
            \begin{block}{BioBERT (Lee et al., 2020)}
                BERT pretrained on biomedical corpora:
                \begin{itemize}
                    \item \textbf{Training Data:}
                    \begin{itemize}
                        \item PubMed abstracts (4.5B words)
                        \item PMC full-text articles (13.5B words)
                    \end{itemize}
                    \item \textbf{Architecture:} BERT-base (110M params)
                    \item \textbf{Tasks:} NER, relation extraction, QA
                \end{itemize}
            \end{block}
            
            \begin{block}{Performance Gains}
                \begin{center}
                    \begin{tabular}{lcc}
                        \toprule
                        \textbf{Task} & \textbf{BERT} & \textbf{BioBERT} \\
                        \midrule
                        NER (F1) & 85.2 & \textcolor{ForestGreen}{89.7} \\
                        RE (F1) & 72.1 & \textcolor{ForestGreen}{77.5} \\
                        QA (F1) & 73.3 & \textcolor{ForestGreen}{81.4} \\
                        \bottomrule
                    \end{tabular}
                \end{center}
            \end{block}
        \end{column}
        \begin{column}{0.5\textwidth}
            \begin{block}{PubMedBERT (Gu et al., 2021)}
                \textbf{From-scratch pretraining} on PubMed:
                \begin{itemize}
                    \item Custom vocabulary from biomedical text
                    \item No transfer from general domain
                    \item Better domain-specific tokenization
                \end{itemize}
            \end{block}
            
            \begin{alertblock}{Key Finding}
                Domain-specific pretraining from scratch outperforms continued pretraining from general models.
            \end{alertblock}
            
            \begin{block}{Our Use Case}
                \begin{itemize}
                    \item Extract medical entities from reports
                    \item Generate clinical terminology embeddings
                    \item Support explanation generation
                \end{itemize}
            \end{block}
        \end{column}
    \end{columns}
\end{frame}

\begin{frame}{Med-PaLM: Medical Large Language Models}
    \begin{columns}
        \begin{column}{0.5\textwidth}
            \begin{block}{Med-PaLM (Singhal et al., 2023)}
                \textbf{Foundation:} PaLM (540B parameters) \\
                \textbf{Fine-tuning:}
                \begin{itemize}
                    \item Instruction tuning on medical QA
                    \item Chain-of-thought prompting
                    \item Human feedback alignment
                \end{itemize}
                
                \vspace{0.5em}
                \textbf{Key Results:}
                \begin{itemize}
                    \item First AI to pass USMLE (>60\%)
                    \item Expert-level on MedQA benchmark
                    \item Clinician-preferred explanations
                \end{itemize}
            \end{block}
        \end{column}
        \begin{column}{0.5\textwidth}
            \begin{block}{Med-PaLM 2 Improvements}
                \begin{itemize}
                    \item 86.5\% on MedQA (vs 67.6\% Med-PaLM)
                    \item Better reasoning chains
                    \item Reduced hallucinations
                    \item Multi-step medical reasoning
                \end{itemize}
            \end{block}
            
            \begin{alertblock}{Explainability Advantage}
                Med-PaLM generates:
                \begin{itemize}
                    \item Clinical rationales for predictions
                    \item Step-by-step diagnostic reasoning
                    \item Patient-friendly explanations
                \end{itemize}
            \end{alertblock}
        \end{column}
    \end{columns}
\end{frame}

\begin{frame}{Vision-Language Models: Bridging Images and Text}
    \begin{block}{The Integration Challenge}
        How do we connect visual features from ViT to language understanding in LLMs?
    \end{block}
    
    \begin{columns}
        \begin{column}{0.5\textwidth}
            \begin{block}{Approach 1: CLIP-Style Alignment}
                \begin{itemize}
                    \item Contrastive learning on image-text pairs
                    \item Shared embedding space
                    \item Zero-shot transfer capabilities
                \end{itemize}
                \textbf{Medical variants:} BiomedCLIP, PubMedCLIP
            \end{block}
            
            \begin{block}{Approach 2: Visual Instruction Tuning}
                \begin{itemize}
                    \item Project visual features into LLM space
                    \item Fine-tune on visual QA tasks
                    \item Examples: LLaVA, MiniGPT-4
                \end{itemize}
            \end{block}
        \end{column}
        \begin{column}{0.5\textwidth}
            \begin{block}{Our Proposed Integration}
                \begin{enumerate}
                    \item Extract features with Swin UNETR
                    \item Project to LLM embedding space
                    \item Fine-tune on medical image-text pairs
                    \item Generate explanations for predictions
                \end{enumerate}
            \end{block}
            
            \begin{alertblock}{Expected Output}
                \textit{``The model detected a 2.3cm enhancing lesion in the right temporal lobe, showing 15\% growth compared to the previous scan. This pattern is consistent with high-grade glioma progression.''}
            \end{alertblock}
        \end{column}
    \end{columns}
\end{frame}

\begin{frame}{LLM Integration Architecture}
    \begin{center}
        \Large
        \textbf{Visual Features $\rightarrow$ Projection Layer $\rightarrow$ LLM $\rightarrow$ Explanation}
    \end{center}
    
    \vspace{1em}
    
    \begin{block}{Mathematical Formulation}
        \begin{center}
            $\mathbf{z}_{\text{visual}} = \text{SwinUNETR}_{\text{encoder}}(\mathbf{x}_{\text{image}})$ \hspace{1em} (Visual Features) \\
            \vspace{0.5em}
            $\mathbf{z}_{\text{proj}} = \mathbf{W}_{\text{proj}} \cdot \mathbf{z}_{\text{visual}} + \mathbf{b}_{\text{proj}}$ \hspace{1em} (Projection) \\
            \vspace{0.5em}
            $\mathbf{y}_{\text{text}} = \text{LLM}([\mathbf{z}_{\text{proj}}; \mathbf{z}_{\text{prompt}}])$ \hspace{1em} (Generation)
        \end{center}
    \end{block}
    
    \begin{columns}
        \begin{column}{0.5\textwidth}
            \begin{block}{Training Strategy}
                \begin{enumerate}
                    \item Freeze ViT encoder (pretrained)
                    \item Train projection layer
                    \item Fine-tune LLM with LoRA
                \end{enumerate}
            \end{block}
        \end{column}
        \begin{column}{0.5\textwidth}
            \begin{block}{Prompt Template}
                \textit{``Given the following tumor features: [VISUAL\_TOKENS], generate a clinical explanation for the observed disease progression.''}
            \end{block}
        \end{column}
    \end{columns}
\end{frame}

% =============================================================================
% TRANSITION 4: Generative Models
% =============================================================================
\begin{frame}[plain]
    \centering
    \vfill
    \begin{beamercolorbox}[center,shadow=true,rounded=true]{title}
        \usebeamerfont{title}
        \textbf{Part IV: Generative Models for Medical Imaging}\\
        \vspace{0.5em}
        \large Diffusion Models and Counterfactual Video Generation
    \end{beamercolorbox}
    \vfill
    
    \vspace{1em}
    \normalsize
    \begin{block}{Why Generative Models?}
        \begin{itemize}
            \item Synthesize realistic disease progression sequences
            \item Enable ``what-if'' counterfactual scenarios
            \item Augment limited longitudinal medical data
        \end{itemize}
    \end{block}
\end{frame}

% =============================================================================
% GENERATIVE MODEL FRAMES
% =============================================================================
\begin{frame}{Denoising Diffusion Probabilistic Models (DDPM)}
    \begin{columns}
        \begin{column}{0.5\textwidth}
            \begin{block}{Core Idea (Ho et al., 2020)}
                Learn to reverse a gradual noising process:
                \begin{enumerate}
                    \item \textbf{Forward:} Gradually add noise to data
                    \item \textbf{Reverse:} Learn to denoise step-by-step
                \end{enumerate}
            \end{block}
            
            \begin{block}{Forward Process}
                \begin{center}
                    $q(\mathbf{x}_t | \mathbf{x}_{t-1}) = \mathcal{N}(\mathbf{x}_t; \sqrt{1-\beta_t}\mathbf{x}_{t-1}, \beta_t\mathbf{I})$
                \end{center}
                \begin{itemize}
                    \item $\beta_t$: Noise schedule (small values)
                    \item After $T$ steps: $\mathbf{x}_T \approx \mathcal{N}(0, \mathbf{I})$
                \end{itemize}
            \end{block}
        \end{column}
        \begin{column}{0.5\textwidth}
            \begin{block}{Reverse Process (Learned)}
                \begin{center}
                    $p_\theta(\mathbf{x}_{t-1} | \mathbf{x}_t) = \mathcal{N}(\mathbf{x}_{t-1}; \boldsymbol{\mu}_\theta(\mathbf{x}_t, t), \sigma_t^2\mathbf{I})$
                \end{center}
                Neural network predicts noise to remove:
                \begin{center}
                    $\boldsymbol{\mu}_\theta(\mathbf{x}_t, t) = \frac{1}{\sqrt{\alpha_t}}\left(\mathbf{x}_t - \frac{\beta_t}{\sqrt{1-\bar{\alpha}_t}}\boldsymbol{\epsilon}_\theta(\mathbf{x}_t, t)\right)$
                \end{center}
            \end{block}
            
            \begin{block}{Training Objective}
                \begin{center}
                    $\mathcal{L} = \mathbb{E}_{t, \mathbf{x}_0, \boldsymbol{\epsilon}}\left[\|\boldsymbol{\epsilon} - \boldsymbol{\epsilon}_\theta(\mathbf{x}_t, t)\|^2\right]$
                \end{center}
                Simple MSE loss on predicted noise.
            \end{block}
        \end{column}
    \end{columns}
\end{frame}

\begin{frame}{Latent Diffusion Models (LDM)}
    \begin{block}{Key Innovation (Rombach et al., 2022)}
        Perform diffusion in compressed latent space instead of pixel space:
        \begin{center}
            $\mathbf{z} = \mathcal{E}(\mathbf{x})$ \hspace{2em} $\hat{\mathbf{x}} = \mathcal{D}(\hat{\mathbf{z}})$
        \end{center}
    \end{block}
    
    \begin{columns}
        \begin{column}{0.5\textwidth}
            \begin{block}{Architecture Components}
                \begin{enumerate}
                    \item \textbf{Autoencoder:} Compress images to latent space
                    \item \textbf{U-Net:} Denoise in latent space
                    \item \textbf{Conditioning:} Text, class, or other modalities
                \end{enumerate}
            \end{block}
            
            \begin{block}{Efficiency Gains}
                \begin{itemize}
                    \item Latent space: $64 \times 64$ vs $512 \times 512$
                    \item 64$\times$ computational reduction
                    \item Enables high-resolution generation
                \end{itemize}
            \end{block}
        \end{column}
        \begin{column}{0.5\textwidth}
            \begin{block}{Cross-Attention Conditioning}
                \begin{center}
                    $\text{Attention}(Q, K, V) = \text{softmax}\left(\frac{QK^T}{\sqrt{d}}\right)V$
                \end{center}
                Where:
                \begin{itemize}
                    \item $Q$: From U-Net features
                    \item $K, V$: From text/condition embeddings
                \end{itemize}
            \end{block}
            
            \begin{alertblock}{Medical Application}
                Condition on:
                \begin{itemize}
                    \item Disease stage (Grade I-IV)
                    \item Treatment type
                    \item Temporal progression
                \end{itemize}
            \end{alertblock}
        \end{column}
    \end{columns}
\end{frame}

\begin{frame}{Video Latent Diffusion Models}
    \begin{block}{Extending LDM to Video (Blattmann et al., 2023)}
        Add temporal modeling to generate coherent video sequences
    \end{block}
    
    \begin{columns}
        \begin{column}{0.5\textwidth}
            \begin{block}{Temporal Architecture}
                \begin{enumerate}
                    \item \textbf{Spatial layers:} Process individual frames
                    \item \textbf{Temporal layers:} Model frame-to-frame dynamics
                    \item \textbf{3D U-Net:} Joint spatial-temporal processing
                \end{enumerate}
            \end{block}
            
            \begin{block}{Temporal Attention}
                \begin{center}
                    $\mathbf{z}_{\text{out}} = \text{TemporalAttn}(\mathbf{z}_1, \mathbf{z}_2, ..., \mathbf{z}_T)$
                \end{center}
                Attention across time dimension ensures:
                \begin{itemize}
                    \item Temporal coherence
                    \item Smooth transitions
                    \item Consistent content
                \end{itemize}
            \end{block}
        \end{column}
        \begin{column}{0.5\textwidth}
            \begin{block}{Disease Progression Modeling}
                \textbf{Inputs:}
                \begin{itemize}
                    \item Baseline scan $\mathbf{x}_0$
                    \item Condition: ``tumor growth over 6 months''
                \end{itemize}
                
                \vspace{0.5em}
                \textbf{Outputs:}
                \begin{itemize}
                    \item Sequence $[\mathbf{x}_0, \mathbf{x}_1, ..., \mathbf{x}_T]$
                    \item Realistic disease progression video
                \end{itemize}
            \end{block}
            
            \begin{alertblock}{Counterfactual Generation}
                \textit{``What would this tumor look like with/without treatment at month 3, 6, 12?''}
            \end{alertblock}
        \end{column}
    \end{columns}
\end{frame}

\begin{frame}{Counterfactual Video Generation: Our Approach}
    \begin{block}{Research Goal}
        Generate realistic ``what-if'' disease progression videos conditioned on treatment scenarios
    \end{block}
    
    \begin{columns}
        \begin{column}{0.5\textwidth}
            \begin{block}{Methodology}
                \begin{enumerate}
                    \item Train Video LDM on longitudinal data
                    \item Condition on treatment embedding
                    \item Generate progression under different scenarios
                    \item Validate with clinical experts
                \end{enumerate}
            \end{block}
            
            \begin{block}{Conditioning Variables}
                \begin{itemize}
                    \item \textcolor{ForestGreen}{Treatment type} (surgery, radiation, chemo)
                    \item \textcolor{violet}{Initial tumor characteristics}
                    \item \textcolor{BrickRed}{Time horizon} (3, 6, 12 months)
                    \item \textcolor{Orange}{Patient demographics}
                \end{itemize}
            \end{block}
        \end{column}
        \begin{column}{0.5\textwidth}
            \begin{block}{Clinical Value}
                \begin{itemize}
                    \item \textbf{Treatment planning:} Visualize expected outcomes
                    \item \textbf{Patient communication:} Show possible futures
                    \item \textbf{Training:} Educate residents on progression patterns
                \end{itemize}
            \end{block}
            
            \begin{alertblock}{Validation Metrics}
                \begin{itemize}
                    \item FID/FVD for visual quality
                    \item Clinical realism assessment
                    \item Progression pattern accuracy
                \end{itemize}
            \end{alertblock}
        \end{column}
    \end{columns}
\end{frame}

% =============================================================================
% TRANSITION 5: Datasets
% =============================================================================
\begin{frame}[plain]
    \centering
    \vfill
    \begin{beamercolorbox}[center,shadow=true,rounded=true]{title}
        \usebeamerfont{title}
        \textbf{Part V: Medical Imaging Datasets}\\
        \vspace{0.5em}
        \large BraTS, TCGA, LIDC-IDRI, and MIMIC-III
    \end{beamercolorbox}
    \vfill
    
    \vspace{1em}
    \normalsize
    \begin{block}{Dataset Requirements}
        \begin{itemize}
            \item High-quality medical images with expert annotations
            \item Longitudinal data for progression modeling
            \item Multi-modal information (imaging + clinical text)
        \end{itemize}
    \end{block}
\end{frame}

% =============================================================================
% DATASET FRAMES
% =============================================================================
\begin{frame}{BraTS Challenge: Brain Tumor Segmentation}
    \begin{columns}
        \begin{column}{0.5\textwidth}
            \begin{block}{Dataset Overview}
                \textbf{Brain Tumor Segmentation Challenge}
                \begin{itemize}
                    \item \textbf{Patients:} 2,000+ cases (2021 version)
                    \item \textbf{Modalities:} T1, T1Gd, T2, T2-FLAIR
                    \item \textbf{Annotations:} Expert neuroradiologists
                    \item \textbf{Resolution:} $240 \times 240 \times 155$
                \end{itemize}
            \end{block}
            
            \begin{block}{Segmentation Labels}
                \begin{itemize}
                    \item \textcolor{ForestGreen}{Whole Tumor (WT)}: Complete tumor extent
                    \item \textcolor{BrickRed}{Tumor Core (TC)}: Non-enhancing + enhancing
                    \item \textcolor{Orange}{Enhancing Tumor (ET)}: Active tumor
                \end{itemize}
            \end{block}
        \end{column}
        \begin{column}{0.5\textwidth}
            \begin{block}{MRI Sequences}
                \begin{center}
                    \begin{tabular}{lp{5cm}}
                        \toprule
                        \textbf{Sequence} & \textbf{Clinical Purpose} \\
                        \midrule
                        T1 & Anatomical structure \\
                        T1Gd & Enhancing tumor (contrast) \\
                        T2 & Edema visualization \\
                        T2-FLAIR & Lesion boundaries \\
                        \bottomrule
                    \end{tabular}
                \end{center}
            \end{block}
            
            \begin{alertblock}{Project Relevance: $\star\star\star\star\star$}
                \textbf{Primary dataset} for:
                \begin{itemize}
                    \item ViT training and evaluation
                    \item Tumor progression modeling
                    \item Benchmark comparisons
                \end{itemize}
            \end{alertblock}
        \end{column}
    \end{columns}
\end{frame}

\begin{frame}{TCGA: The Cancer Genome Atlas}
    \begin{columns}
        \begin{column}{0.5\textwidth}
            \begin{block}{Dataset Overview}
                \textbf{Multi-Omics Cancer Database}
                \begin{itemize}
                    \item \textbf{Patients:} 11,000+ cases
                    \item \textbf{Cancer Types:} 33 different types
                    \item \textbf{Data Types:}
                    \begin{itemize}
                        \item Pathology images (WSI)
                        \item Genomic data
                        \item Clinical records
                        \item Radiology images
                    \end{itemize}
                \end{itemize}
            \end{block}
            
            \begin{block}{Relevant Subsets}
                \begin{itemize}
                    \item \textbf{TCGA-GBM:} Glioblastoma (600+ cases)
                    \item \textbf{TCGA-LGG:} Low-grade glioma (500+ cases)
                    \item \textbf{TCGA-LUAD:} Lung adenocarcinoma
                \end{itemize}
            \end{block}
        \end{column}
        \begin{column}{0.5\textwidth}
            \begin{block}{Multi-Modal Integration}
                \begin{center}
                    \begin{tabular}{lc}
                        \toprule
                        \textbf{Modality} & \textbf{Availability} \\
                        \midrule
                        Pathology WSI & $\checkmark$ \\
                        MRI/CT Images & $\checkmark$ \\
                        Genomic Data & $\checkmark$ \\
                        Clinical Reports & $\checkmark$ \\
                        Survival Data & $\checkmark$ \\
                        \bottomrule
                    \end{tabular}
                \end{center}
            \end{block}
            
            \begin{alertblock}{Project Relevance: $\star\star\star\star$}
                \textbf{Secondary dataset} for:
                \begin{itemize}
                    \item Multi-modal model training
                    \item Cross-cancer generalization
                    \item Clinical text integration
                \end{itemize}
            \end{alertblock}
        \end{column}
    \end{columns}
\end{frame}

\begin{frame}{LIDC-IDRI: Lung Image Database Consortium}
    \begin{columns}
        \begin{column}{0.5\textwidth}
            \begin{block}{Dataset Overview}
                \textbf{Lung CT Nodule Database}
                \begin{itemize}
                    \item \textbf{Patients:} 1,018 cases
                    \item \textbf{Nodules:} 2,669 annotated lesions
                    \item \textbf{Annotators:} 4 thoracic radiologists
                    \item \textbf{Slice Thickness:} 0.6 - 5.0 mm
                \end{itemize}
            \end{block}
            
            \begin{block}{Annotation Details}
                For each nodule $\geq$ 3mm:
                \begin{itemize}
                    \item 3D segmentation boundaries
                    \item Malignancy rating (1-5)
                    \item Texture, margin, subtlety scores
                    \item Likelihood of malignancy
                \end{itemize}
            \end{block}
        \end{column}
        \begin{column}{0.5\textwidth}
            \begin{block}{Multi-Annotator Consensus}
                \begin{itemize}
                    \item 4 independent annotations per nodule
                    \item Inter-observer variability studied
                    \item Consensus for training/evaluation
                \end{itemize}
            \end{block}
            
            \begin{block}{Malignancy Distribution}
                \begin{center}
                    \begin{tabular}{cc}
                        \toprule
                        \textbf{Score} & \textbf{Percentage} \\
                        \midrule
                        1 (Benign) & 23\% \\
                        2-3 & 35\% \\
                        4-5 (Malignant) & 42\% \\
                        \bottomrule
                    \end{tabular}
                \end{center}
            \end{block}
            
            \begin{alertblock}{Project Relevance: $\star\star\star$}
                \textbf{Validation dataset} for model generalization
            \end{alertblock}
        \end{column}
    \end{columns}
\end{frame}

\begin{frame}{MIMIC-III: Clinical Text Database}
    \begin{columns}
        \begin{column}{0.5\textwidth}
            \begin{block}{Dataset Overview}
                \textbf{Medical Information Mart for Intensive Care}
                \begin{itemize}
                    \item \textbf{Patients:} 46,520 ICU admissions
                    \item \textbf{Time Span:} 2001-2012
                    \item \textbf{Hospital:} Beth Israel Deaconess
                \end{itemize}
            \end{block}
            
            \begin{block}{Data Types}
                \begin{itemize}
                    \item \textbf{Clinical Notes:} 2M+ documents
                    \item \textbf{Radiology Reports:} 500K+ reports
                    \item \textbf{Discharge Summaries}
                    \item \textbf{Lab Results}
                    \item \textbf{Medications}
                \end{itemize}
            \end{block}
        \end{column}
        \begin{column}{0.5\textwidth}
            \begin{block}{Text Categories}
                \begin{center}
                    \begin{tabular}{lc}
                        \toprule
                        \textbf{Note Type} & \textbf{Count} \\
                        \midrule
                        Nursing Notes & 800K+ \\
                        Radiology & 500K+ \\
                        Physician Notes & 400K+ \\
                        Discharge & 50K+ \\
                        \bottomrule
                    \end{tabular}
                \end{center}
            \end{block}
            
            \begin{alertblock}{Project Relevance: $\star\star\star\star$}
                \textbf{Critical for LLM training:}
                \begin{itemize}
                    \item Clinical language patterns
                    \item Radiology report templates
                    \item Explanation generation training
                \end{itemize}
            \end{alertblock}
        \end{column}
    \end{columns}
\end{frame}

\begin{frame}{Dataset Summary and Selection}
    \begin{center}
        \begin{tabular}{@{}lcccc@{}}
            \toprule
            \textbf{Dataset} & \textbf{Modality} & \textbf{Size} & \textbf{Use Case} & \textbf{Relevance} \\
            \midrule
            BraTS & Brain MRI & 2,000+ & Segmentation, Progression & $\star\star\star\star\star$ \\
            TCGA & Multi-modal & 11,000+ & Multi-cancer analysis & $\star\star\star\star$ \\
            LIDC-IDRI & Lung CT & 1,018 & Generalization testing & $\star\star\star$ \\
            MIMIC-III & Clinical text & 46,520 & LLM training & $\star\star\star\star$ \\
            \bottomrule
        \end{tabular}
    \end{center}
    
    \vspace{1em}
    
    \begin{columns}
        \begin{column}{0.5\textwidth}
            \begin{block}{Primary Focus}
                \begin{itemize}
                    \item \textbf{BraTS:} Main imaging dataset
                    \item \textbf{MIMIC-III:} Clinical text source
                \end{itemize}
            \end{block}
        \end{column}
        \begin{column}{0.5\textwidth}
            \begin{block}{Secondary/Validation}
                \begin{itemize}
                    \item \textbf{TCGA:} Multi-modal extension
                    \item \textbf{LIDC-IDRI:} Cross-domain validation
                \end{itemize}
            \end{block}
        \end{column}
    \end{columns}
\end{frame}

% =============================================================================
% TRANSITION 6: Project Plan
% =============================================================================
\begin{frame}[plain]
    \centering
    \vfill
    \begin{beamercolorbox}[center,shadow=true,rounded=true]{title}
        \usebeamerfont{title}
        \textbf{Part VI: Research Project Plan}\\
        \vspace{0.5em}
        \large Objectives, Timeline, and Expected Deliverables
    \end{beamercolorbox}
    \vfill
    
    \vspace{1em}
    \normalsize
    \textit{20-week Mitacs Globalink Research Internship}\\
    April 1 - August 19, 2026
\end{frame}

% =============================================================================
% PROJECT PLAN FRAMES
% =============================================================================
\begin{frame}{Research Objectives Overview}
    \begin{block}{Main Goal}
        Develop an explainable AI system for disease progression modeling and counterfactual video generation using vision-language models
    \end{block}
    
    \begin{center}
        \begin{tabular}{@{}clc@{}}
            \toprule
            \textbf{Obj.} & \textbf{Description} & \textbf{Duration} \\
            \midrule
            1 & Data Pipeline Development & Weeks 1-4 \\
            2 & ViT-Based Feature Extraction & Weeks 5-8 \\
            3 & LLM Integration for Explainability & Weeks 9-12 \\
            4 & Counterfactual Video Generation & Weeks 13-16 \\
            5 & Evaluation and Documentation & Weeks 17-20 \\
            \bottomrule
        \end{tabular}
    \end{center}
\end{frame}

\begin{frame}{Objective 1: Data Pipeline Development (Weeks 1-4)}
    \begin{columns}
        \begin{column}{0.5\textwidth}
            \begin{block}{Tasks}
                \begin{enumerate}
                    \item Download and organize BraTS dataset
                    \item Implement preprocessing pipeline:
                    \begin{itemize}
                        \item Skull stripping
                        \item Intensity normalization
                        \item Registration to template
                    \end{itemize}
                    \item Create data loaders for PyTorch
                    \item Set up experiment tracking (W\&B)
                \end{enumerate}
            \end{block}
            
            \begin{block}{Key Technologies}
                \begin{itemize}
                    \item MONAI for medical imaging
                    \item PyTorch DataLoader
                    \item NiBabel for NIfTI handling
                \end{itemize}
            \end{block}
        \end{column}
        \begin{column}{0.5\textwidth}
            \begin{block}{Deliverables}
                \begin{itemize}
                    \item $\checkmark$ Organized dataset structure
                    \item $\checkmark$ Preprocessing scripts
                    \item $\checkmark$ Data augmentation pipeline
                    \item $\checkmark$ Train/val/test splits
                \end{itemize}
            \end{block}
            
            \begin{alertblock}{Success Criteria}
                \begin{itemize}
                    \item All data accessible via unified API
                    \item Preprocessing validated visually
                    \item Reproducible data loading
                \end{itemize}
            \end{alertblock}
        \end{column}
    \end{columns}
\end{frame}

\begin{frame}{Objective 2: ViT-Based Feature Extraction (Weeks 5-8)}
    \begin{columns}
        \begin{column}{0.5\textwidth}
            \begin{block}{Tasks}
                \begin{enumerate}
                    \item Implement Swin UNETR architecture
                    \item Load pretrained weights
                    \item Fine-tune on BraTS dataset
                    \item Extract intermediate representations
                    \item Analyze attention patterns
                \end{enumerate}
            \end{block}
            
            \begin{block}{Architecture Details}
                \begin{itemize}
                    \item \textbf{Encoder:} Swin Transformer
                    \item \textbf{Decoder:} CNN upsampling
                    \item \textbf{Output:} Multi-class segmentation
                \end{itemize}
            \end{block}
        \end{column}
        \begin{column}{0.5\textwidth}
            \begin{block}{Target Metrics}
                \begin{center}
                    \begin{tabular}{lc}
                        \toprule
                        \textbf{Region} & \textbf{Target Dice} \\
                        \midrule
                        Whole Tumor & $\geq$ 0.90 \\
                        Tumor Core & $\geq$ 0.85 \\
                        Enhancing & $\geq$ 0.80 \\
                        \bottomrule
                    \end{tabular}
                \end{center}
            \end{block}
            
            \begin{alertblock}{Deliverables}
                \begin{itemize}
                    \item Trained Swin UNETR model
                    \item Feature extraction module
                    \item Attention visualization tools
                \end{itemize}
            \end{alertblock}
        \end{column}
    \end{columns}
\end{frame}

\begin{frame}{Objective 3: LLM Integration for Explainability (Weeks 9-12)}
    \begin{columns}
        \begin{column}{0.5\textwidth}
            \begin{block}{Tasks}
                \begin{enumerate}
                    \item Design vision-language projection layer
                    \item Fine-tune LLM on medical explanations
                    \item Create prompt templates for reports
                    \item Implement explanation generation pipeline
                    \item Evaluate explanation quality
                \end{enumerate}
            \end{block}
            
            \begin{block}{LLM Options}
                \begin{itemize}
                    \item \textbf{Primary:} LLaMA-2 + LoRA fine-tuning
                    \item \textbf{Alternative:} Mistral-7B
                    \item \textbf{Training Data:} MIMIC-III reports
                \end{itemize}
            \end{block}
        \end{column}
        \begin{column}{0.5\textwidth}
            \begin{block}{Evaluation Metrics}
                \begin{itemize}
                    \item \textbf{BLEU/ROUGE:} Text quality
                    \item \textbf{Clinical Accuracy:} Expert review
                    \item \textbf{Faithfulness:} Explanation-prediction alignment
                \end{itemize}
            \end{block}
            
            \begin{alertblock}{Expected Output}
                \textit{``Analysis shows a 23mm enhancing mass in the right temporal lobe with surrounding edema. Compared to baseline, tumor volume increased 15\%, suggesting disease progression. Key features: irregular margins, heterogeneous enhancement.''}
            \end{alertblock}
        \end{column}
    \end{columns}
\end{frame}

\begin{frame}{Objective 4: Counterfactual Video Generation (Weeks 13-16)}
    \begin{columns}
        \begin{column}{0.5\textwidth}
            \begin{block}{Tasks}
                \begin{enumerate}
                    \item Implement Video Latent Diffusion Model
                    \item Train on longitudinal tumor data
                    \item Design conditioning mechanism
                    \item Generate counterfactual sequences
                    \item Validate clinical realism
                \end{enumerate}
            \end{block}
            
            \begin{block}{Conditioning Variables}
                \begin{itemize}
                    \item Treatment type (surgery, radiation, chemo)
                    \item Time horizon (3, 6, 12 months)
                    \item Baseline tumor characteristics
                \end{itemize}
            \end{block}
        \end{column}
        \begin{column}{0.5\textwidth}
            \begin{block}{Technical Approach}
                \begin{enumerate}
                    \item 3D VAE for latent compression
                    \item Temporal U-Net for diffusion
                    \item Cross-attention for conditioning
                    \item Classifier-free guidance
                \end{enumerate}
            \end{block}
            
            \begin{alertblock}{Validation}
                \begin{itemize}
                    \item FVD score for video quality
                    \item Clinical expert assessment
                    \item Progression pattern accuracy
                \end{itemize}
            \end{alertblock}
        \end{column}
    \end{columns}
\end{frame}

\begin{frame}{Objective 5: Evaluation and Documentation (Weeks 17-20)}
    \begin{columns}
        \begin{column}{0.5\textwidth}
            \begin{block}{Evaluation Tasks}
                \begin{enumerate}
                    \item Comprehensive model testing
                    \item Ablation studies
                    \item Cross-dataset generalization
                    \item Clinical expert evaluation
                    \item User study (if time permits)
                \end{enumerate}
            \end{block}
            
            \begin{block}{Documentation}
                \begin{itemize}
                    \item Technical report
                    \item Code documentation
                    \item Model cards
                    \item Usage tutorials
                \end{itemize}
            \end{block}
        \end{column}
        \begin{column}{0.5\textwidth}
            \begin{block}{Expected Deliverables}
                \begin{itemize}
                    \item $\checkmark$ Trained models (ViT, LLM, Video LDM)
                    \item $\checkmark$ Evaluation results
                    \item $\checkmark$ GitHub repository
                    \item $\checkmark$ Technical report / draft paper
                    \item $\checkmark$ Final presentation
                \end{itemize}
            \end{block}
            
            \begin{alertblock}{Publication Target}
                \begin{itemize}
                    \item MICCAI 2026 (Medical Imaging)
                    \item CVPR 2027 (Computer Vision)
                    \item Nature Medicine (if results strong)
                \end{itemize}
            \end{alertblock}
        \end{column}
    \end{columns}
\end{frame}

\begin{frame}{Project Timeline: Gantt Chart Overview}
    \begin{center}
        \begin{tabular}{|l|c|c|c|c|c|}
            \hline
            \textbf{Phase} & \textbf{W1-4} & \textbf{W5-8} & \textbf{W9-12} & \textbf{W13-16} & \textbf{W17-20} \\
            \hline
            Data Pipeline & \cellcolor{ForestGreen!60} & & & & \\
            \hline
            ViT Training & & \cellcolor{ForestGreen!60} & & & \\
            \hline
            LLM Integration & & & \cellcolor{ForestGreen!60} & & \\
            \hline
            Video Generation & & & & \cellcolor{ForestGreen!60} & \\
            \hline
            Evaluation & & & & & \cellcolor{ForestGreen!60} \\
            \hline
            Documentation & \cellcolor{blue!20} & \cellcolor{blue!20} & \cellcolor{blue!20} & \cellcolor{blue!20} & \cellcolor{blue!60} \\
            \hline
            Meetings & \cellcolor{orange!30} & \cellcolor{orange!30} & \cellcolor{orange!30} & \cellcolor{orange!30} & \cellcolor{orange!30} \\
            \hline
        \end{tabular}
    \end{center}
    
    \vspace{1em}
    
    \begin{columns}
        \begin{column}{0.5\textwidth}
            \begin{block}{Key Milestones}
                \begin{itemize}
                    \item Week 4: Data pipeline complete
                    \item Week 8: ViT baseline trained
                    \item Week 12: LLM integration done
                    \item Week 16: Video generation working
                \end{itemize}
            \end{block}
        \end{column}
        \begin{column}{0.5\textwidth}
            \begin{block}{Risk Mitigation}
                \begin{itemize}
                    \item Weekly supervisor meetings
                    \item Incremental development
                    \item Fallback to simpler models if needed
                \end{itemize}
            \end{block}
        \end{column}
    \end{columns}
\end{frame}

% =============================================================================
% REFERENCES SUMMARY
% =============================================================================
\begin{frame}{Key References: Helpfulness Scores}
    \begin{center}
        \footnotesize
        \begin{tabular}{@{}p{4.5cm}p{3cm}cc@{}}
            \toprule
            \textbf{Paper} & \textbf{Topic} & \textbf{Score} & \textbf{Objectives} \\
            \midrule
            Dosovitskiy et al., 2020 & Vision Transformer & 9/10 & O2 \\
            Liu et al., 2021 & Swin Transformer & 10/10 & O2 \\
            Hatamizadeh et al., 2022 & UNETR & 10/10 & O2 \\
            Tang et al., 2022 & Swin UNETR & 10/10 & O2 \\
            Lee et al., 2020 & BioBERT & 8/10 & O3 \\
            Singhal et al., 2023 & Med-PaLM & 9/10 & O3 \\
            Ho et al., 2020 & DDPM & 9/10 & O4 \\
            Rombach et al., 2022 & Latent Diffusion & 10/10 & O4 \\
            Blattmann et al., 2023 & Video LDM & 10/10 & O4 \\
            \bottomrule
        \end{tabular}
    \end{center}
    
    \vspace{0.5em}
    \footnotesize
    \textbf{Legend:} O2 = ViT Features, O3 = LLM Integration, O4 = Video Generation
\end{frame}

% =============================================================================
% CONCLUSION
% =============================================================================
\begin{frame}{Conclusion: Project Summary}
    \begin{block}{Key Contributions}
        \begin{enumerate}
            \item \textbf{Novel Framework:} Integration of ViT + LLM + Diffusion for medical imaging
            \item \textbf{Explainability:} Natural language explanations for visual predictions
            \item \textbf{Counterfactual Generation:} ``What-if'' disease progression videos
            \item \textbf{Clinical Utility:} Tools for treatment planning and patient communication
        \end{enumerate}
    \end{block}
    
    \begin{columns}
        \begin{column}{0.5\textwidth}
            \begin{block}{Expected Impact}
                \begin{itemize}
                    \item Bridge AI-clinician trust gap
                    \item Enable personalized medicine
                    \item Advance medical AI interpretability
                \end{itemize}
            \end{block}
        \end{column}
        \begin{column}{0.5\textwidth}
            \begin{alertblock}{Next Steps}
                \begin{itemize}
                    \item Begin Phase 1 implementation
                    \item Set up computing environment
                    \item Download and preprocess BraTS
                \end{itemize}
            \end{alertblock}
        \end{column}
    \end{columns}
\end{frame}

% =============================================================================
% REFERENCES
% =============================================================================
\begin{frame}[allowframebreaks]{References}
\frametitle{References}
\begin{thebibliography}{20}
\setlength{\itemsep}{0.3em}
\footnotesize

\bibitem{dosovitskiy2020vit}
Alexey Dosovitskiy et al.
\newblock An Image is Worth 16x16 Words: Transformers for Image Recognition at Scale.
\newblock \emph{ICLR}, 2021.

\bibitem{liu2021swin}
Ze Liu et al.
\newblock Swin Transformer: Hierarchical Vision Transformer using Shifted Windows.
\newblock \emph{ICCV}, 2021.

\bibitem{hatamizadeh2022unetr}
Ali Hatamizadeh et al.
\newblock UNETR: Transformers for 3D Medical Image Segmentation.
\newblock \emph{WACV}, 2022.

\bibitem{tang2022swinunetr}
Yucheng Tang et al.
\newblock Self-Supervised Pre-Training of Swin Transformers for 3D Medical Image Analysis.
\newblock \emph{CVPR}, 2022.

\bibitem{lee2020biobert}
Jinhyuk Lee et al.
\newblock BioBERT: A Pre-trained Biomedical Language Representation Model.
\newblock \emph{Bioinformatics}, 2020.

\bibitem{singhal2023medpalm}
Karan Singhal et al.
\newblock Large Language Models Encode Clinical Knowledge.
\newblock \emph{Nature}, 2023.

\bibitem{ho2020ddpm}
Jonathan Ho et al.
\newblock Denoising Diffusion Probabilistic Models.
\newblock \emph{NeurIPS}, 2020.

\bibitem{rombach2022ldm}
Robin Rombach et al.
\newblock High-Resolution Image Synthesis with Latent Diffusion Models.
\newblock \emph{CVPR}, 2022.

\bibitem{blattmann2023videoldm}
Andreas Blattmann et al.
\newblock Align your Latents: High-Resolution Video Synthesis with Latent Diffusion Models.
\newblock \emph{CVPR}, 2023.

\end{thebibliography}
\end{frame}

% =============================================================================
% THANK YOU
% =============================================================================
\begin{frame}[standout]
    \centering
    \Huge Thank You! \\
    \vspace{1cm}
    \Large Questions \& Discussion
    
    \vspace{2cm}
    \normalsize
    Mohamed Amine \\
    TÉLUQ University, Montreal \\
    Mitacs Globalink Research Internship 2026
\end{frame}

\end{document}
